\documentclass[11pt]{article}
\usepackage[margin=1in]{geometry}
\usepackage{amsmath,amssymb,amsthm,bm}
\usepackage{graphicx}
\usepackage{booktabs}
\usepackage{multirow}
\usepackage{hyperref}
\usepackage{enumitem}
\usepackage{algorithm}
\usepackage{algpseudocode}

\newcommand{\C}{\mathbb{C}}
\newcommand{\R}{\mathbb{R}}
\newcommand{\E}{\mathbb{E}}
\newcommand{\Hc}{\mathbf{H}}
\newcommand{\yv}{\mathbf{y}}
\newcommand{\xv}{\mathbf{x}}
\newcommand{\zv}{\mathbf{z}}
\newcommand{\av}{\mathbf{a}}
\newcommand{\hv}{\mathbf{h}}
\newcommand{\uv}{\mathbf{u}}
\newcommand{\bs}{\mathbf{b}}
\newcommand{\softplus}{\operatorname{softplus}}
\DeclareMathOperator{\NMSE}{NMSE}
\DeclareMathOperator{\FFT}{FFT}
\DeclareMathOperator{\IFFT}{IFFT}

\title{\textbf{Beamspace Selective State-Space World Model for\\ Wireless Channel Estimation and Prediction}\\[0.4em]
\large End-to-End Methodology}
\author{}
\date{}

\begin{document}
\maketitle

\begin{abstract}
We develop a \emph{world model} for MIMO-OFDM channels: an action-conditioned model of channel
dynamics whose one-step prediction serves as a temporal prior for estimation. The model operates in
the \textbf{beamspace} (angle--delay) domain, where the channel is sparse and a compact latent is
reconstruction-faithful; a Mamba-style \emph{selective state-space} (SSM) backbone rolls the latent
forward under the receiver-motion action and \emph{predicts the future channel}, and a
noise-conditioned convolutional head fuses that prediction with the noisy observation to estimate the
current channel. Trained end-to-end across $8\times$ A100 GPUs on
$6.4{\times}10^4$ ray-traced Sionna~RT channel sequences over six city scenes, the estimator
attains lower NMSE than the optimal linear (MMSE) estimator \emph{at every SNR from $-5$ to $30$~dB},
\emph{in-distribution (over 3 seeds) and on every one of six held-out (unseen) scenes} (leave-one-
scene-out, $42/42$ wins), without access to the channel covariance or noise power that MMSE requires;
the predictor likewise beats persistence and linear AR(1) at every horizon $k=1,\dots,6$. \S\ref{sec:discarded} records, honestly, the components
we tried and discarded on the evidence (a JEPA joint-embedding objective and a lossy pretrained
encoder), and the controls (a same-capacity supervised baseline; an out-of-distribution split) that
drove the final design.
\end{abstract}

\tableofcontents

%==================================================================
\section{Problem formulation}\label{sec:problem}

\subsection{Signal model}
A single-cell MIMO-OFDM link has a base station of $N_a$ antennas and a single-antenna user
equipment. The frequency-domain channel at time $t$ is
\begin{equation}
\Hc_t \in \C^{N_a \times N_s},\qquad
[\Hc_t]_{m,k} = \sum_{\ell=1}^{L} \alpha_{\ell}\,
e^{\,j 2\pi m\phi_\ell/N_a}\,e^{-j 2\pi k\tau_\ell/N_s},
\end{equation}
with $L$ paths of complex gain, angle, and delay $(\alpha_\ell,\phi_\ell,\tau_\ell)$. We stack
real/imag parts, $\mathbf{o}_t=[\Re\Hc_t;\Im\Hc_t]\in\R^{2\times N_a\times N_s}$. The pilot
observation is $\yv_t=\Hc_t+\mathbf{n}_t$, $\mathbf{n}_t\!\sim\!\mathcal{CN}(0,\sigma^2\mathbf I)$,
with $\mathrm{SNR}=10\log_{10}(\bar P/\sigma^2)$, $\bar P=\E\|\Hc_t\|_F^2/(N_aN_s)$.

\subsection{Actions and temporal structure}
The receiver moves; consecutive channels are correlated and their evolution is driven by the motion.
The \textbf{action} is the velocity descriptor $\av_t=[v_x,v_y,\|\mathbf v\|,\theta]\in\R^{d_a}$,
$d_a=4$. A trajectory yields $\{(\mathbf{o}_t,\av_t)\}_{t=0}^{T-1}$.

\subsection{Two tasks and metric}
\textbf{Estimation:} recover clean $\Hc_t$ from noisy $\yv_t$. \textbf{Prediction:} forecast the
channel $k$ steps ahead from history and planned actions, without observing the future. Metric
throughout: $\NMSE(\hat\Hc)=\E\|\hat\Hc-\Hc\|_F^2/\E\|\Hc\|_F^2$.

%==================================================================
\section{Architecture overview}\label{sec:arch}

The world model is one action-conditioned dynamics model that serves both tasks:
\begin{equation}
\mathbf{o}_t\!\xrightarrow{\text{2D FFT}}\!\bs_t\!\xrightarrow{f_\theta}\!\xv_t
\!\xrightarrow{\text{selective SSM}}\!(\zv_t,\hv_t)
\!\xrightarrow{\text{roll }k\text{ w/ }\av}\!\hat{\bs}_{t+k}\!\xrightarrow{g_\omega}\!\text{predicted channel},
\end{equation}
and estimation is \emph{predict-then-correct}: the one-step prediction $\hat{\bs}_t$ from history is
fused with the noisy beamspace observation by a learned gain to produce $\hat\Hc_t$. Every module is
load-bearing for both tasks --- the encoder/SSM/predictor produce the prior, the fusion head corrects
it --- over a single sparse, faithful representation.

%==================================================================
\section{Beamspace representation}\label{sec:beamspace}

The model operates in \textbf{beamspace} rather than the raw antenna--subcarrier grid. For
$\Hc\in\C^{N_a\times N_s}$ the orthonormal 2-D DFT
\begin{equation}
\bs = \FFT_{\text{sub}}\!\big(\IFFT_{\text{ant}}(\Hc)\big),\qquad
\Hc = \FFT_{\text{ant}}\!\big(\IFFT_{\text{sub}}(\bs)\big)
\end{equation}
(IFFT over antennas $\to$ angle/beam; FFT over subcarriers $\to$ delay) maps paths to a few
$(\text{angle},\text{delay})$ taps. It is orthonormal (energy-preserving, exactly invertible). On the
Sionna channels this is strongly sparse and reconstruction-faithful (Table~\ref{tab:beam}): $90\%$ of
the energy lies in $5.4\%$ of taps, and a $5\%$-tap approximation reaches $0.0024$ NMSE. Three
consequences make this the right domain: (i) a \emph{compact} latent can carry the full channel, so
estimation and prediction need not compete for capacity; (ii) receiver motion becomes a per-tap
\emph{linear phase progression} --- scene-invariant physics that supports out-of-distribution
generalization; (iii) a prediction made in beamspace is directly a good channel estimate, enabling
predict-then-correct.

\begin{table}[h]
\centering
\begin{tabular}{lcccc}
\toprule
kept taps & $2\%$ & $5\%$ & $10\%$ & round-trip \\
\midrule
reconstruction NMSE & $0.034$ & $\mathbf{0.0024}$ & $0.0010$ & $\sim\!10^{-14}$ \\
\bottomrule
\end{tabular}
\caption{Beamspace sparsity/faithfulness on real Sionna channels ($90\%$ energy in $5.4\%$ of taps).}
\label{tab:beam}
\end{table}

%==================================================================
\section{Module 1: Encoder / decoder}\label{sec:encoder}

The encoder is a small convolutional network over the beamspace grid,
$f_\theta:\R^{2\times N_a\times N_s}\!\to\R^{d}$ (two conv blocks, spatial pooling, linear
projection). A symmetric decoder $g_\omega:\R^{d}\!\to\R^{2\times N_a\times N_s}$ reconstructs the
beamspace channel from a latent. Both are trained from scratch; because beamspace is faithful, the
latent is reconstruction-capable (unlike the pretrained image-domain encoder of \S\ref{sec:discarded},
which was lossy).

%==================================================================
\section{Module 2: SelectionNet (input-dependent SSM parameters)}\label{sec:selection}

The dynamics are made \emph{selective} (Mamba-style) by conditioning the diagonal state-space
parameters on the action. SelectionNet $h_\eta:\R^{d_a}\!\to(\R^{n})^4$ (MLP trunk, four heads)
outputs, for state dimension $n$,
\begin{align}
\mathbf{A}_t &= -\softplus(\cdot)\prec\mathbf 0 \ \ \text{(stable poles)}, &
\bm{\Delta}_t &= \softplus(\cdot)\succ\mathbf 0 \ \ \text{(positive step)}, \\
\mathbf{B}_t &= \mathrm{Lin}(\cdot), & \mathbf{C}_t &= \mathrm{Lin}(\cdot).
\end{align}
Negativity of $\mathbf A_t$ guarantees a stable discretization for any action; positivity of
$\bm\Delta_t$ a valid step. HiPPO/S4-style initialization spreads the timescales: the $\mathbf A$-bias
targets $\{-1,\dots,-n\}$ and the $\bm\Delta$-bias is the inverse-softplus of a log-uniform sample in
$[10^{-3},10^{-1}]$; head weights are small but nonzero so all four parameters stay action-dependent.

%==================================================================
\section{Module 3: Selective SSM (temporal backbone)}\label{sec:ssm}

Each of the $n$ diagonal state channels is a scalar system discretized by zero-order hold with step
$\Delta_t$:
\begin{equation}
\bar A_t=\exp(\Delta_t A_t)\in(0,1),\qquad \bar B_t=(\bar A_t-1)A_t^{-1}B_t,
\end{equation}
stable since $A_t<0$. With input mixing $\uv_t=\mathrm{Lin}([\xv_t;\av_t])$,
\begin{equation}
\hv_t=\bar A_t\odot\hv_{t-1}+\bar B_t\odot\uv_t,\qquad
\zv_t=\mathrm{Lin}\big(\mathrm{LN}(\mathbf C_t\odot\hv_t+\mathbf D\odot\uv_t)\big).
\end{equation}
The scan is causal, selective, and stable over long horizons; the recurrent state $\hv_t$ is exposed
so the predictor can warm-start from real history rather than a single latent.

%==================================================================
\section{Module 4: Predictor (future-channel imagination)}\label{sec:predictor}

The predictor rolls the latent $k$ steps forward \emph{using planned actions only} (the future is
unobserved), warm-started from the SSM state $\hv_t$. For $j=0,\dots,k-1$ with
$(\mathbf A_j,\mathbf B_j,\mathbf C_j,\Delta_j)=h_\eta(\av_{t+j})$ (the \emph{shared} SelectionNet):
\begin{equation}
\hv_{j+1}=\bar A_j\odot\hv_j+\bar B_j\odot\mathrm{Lin}(\av_{t+j}),\qquad
\hat{\bs}_{t+k}=g_\omega\big(\mathrm{Lin}(\mathbf y_k)\big),
\end{equation}
decoding to the \emph{future channel in beamspace}. This is a reconstruction-based world model
(predict the future \emph{state}), not a joint-embedding one (\S\ref{sec:discarded}).

\paragraph{Multi-horizon training.}
Training at a single fixed horizon overfits to it and degrades beyond it (a $k{=}3$-only predictor at
$k{=}6$ reached $0.213$ NMSE, worse than persistence). We sample $k\sim\mathcal U\{1,\dots,6\}$ per
step. Evaluated in \emph{channel space} on the 6-scene stress dataset, the learned predictor beats
both persistence and a channel-space linear AR(1) at \emph{every} horizon $k=1,\dots,6$, and the
margin \emph{widens} with $k$ --- the world-model signature: persistence and AR(1) both decay with
horizon while the SSM prediction stays nearly flat (Table~\ref{tab:horizon}). Notably AR(1) crosses
\emph{above} persistence by $k{=}5$ (linear extrapolation compounds error), whereas the nonlinear
action-driven roll does not.

\begin{table}[h]
\centering
\begin{tabular}{lcccccc}
\toprule
horizon $k$ & 1 & 2 & 3 & 4 & 5 & 6 \\
\midrule
persistence & $0.374$ & $0.409$ & $0.431$ & $0.452$ & $0.479$ & $0.509$ \\
AR(1)       & $0.295$ & $0.342$ & $0.395$ & $0.457$ & $0.517$ & $0.573$ \\
\textbf{predictor} & $\mathbf{0.249}$ & $\mathbf{0.260}$ & $\mathbf{0.272}$ & $\mathbf{0.284}$ & $\mathbf{0.296}$ & $\mathbf{0.310}$ \\
\bottomrule
\end{tabular}
\caption{Held-out channel-space prediction NMSE vs.\ horizon (lower better), in-distribution.
Beats both baselines at every horizon; margin widens with $k$.}
\label{tab:horizon}
\end{table}

%==================================================================
\section{Module 5: Predict-then-correct estimation head}\label{sec:heads}

Estimation fuses the world-model prior with the observation, all in beamspace. Given the noisy
beamspace observation $\bs^y=\mathrm{beam}(\yv)$, the one-step prior $\hat{\bs}_t$ predicted from
clean history (\S\ref{sec:predictor}), and per-sample noise variance $\sigma^2$, a learned per-element
\emph{Kalman gain} $\mathbf g\in[0,1]^{2\times N_a\times N_s}$ combines them, followed by a small conv
refinement:
\begin{equation}
\mathbf g=\text{Conv}\big([\bs^y;\hat{\bs}_t;\sigma^2\mathbf 1]\big),\quad
\hat{\bs}=\mathbf g\odot\bs^y+(1-\mathbf g)\odot\hat{\bs}_t+\text{Conv}(\cdot),\quad
\hat\Hc=\mathrm{beam}^{-1}(\hat{\bs}).
\end{equation}
The refinement conv is zero-initialized so the head starts at the gated fusion. The gate is conditioned
on $\sigma^2$ (which MMSE also receives), so it trusts the observation at high SNR and the prior at low
SNR. Crucially, the prior is temporal information a single-snapshot estimator (or MMSE) cannot access.

%==================================================================
\section{Classical baselines}\label{sec:baselines}

\paragraph{Least squares.} With unit pilots $\hat\Hc_{\text{LS}}=\yv$, $\NMSE_{\text{LS}}=10^{-\mathrm{SNR}/10}$.

\paragraph{Linear MMSE (Wiener).} With $\hv=\mathrm{vec}(\Hc)$, covariance $\mathbf R=\E[\hv\hv^H]$
(estimated from training channels) and noise $\sigma^2$,
\begin{equation}
\hat\hv_{\text{MMSE}}=\mathbf R(\mathbf R+\sigma^2\mathbf I)^{-1}\yv_{\text{vec}}.
\end{equation}
MMSE is the optimal \emph{linear} estimator and is handed both $\mathbf R$ and $\sigma^2$; our model
receives only the scalar $\sigma^2$ and no covariance.

\paragraph{Supervised control.} An \emph{identical} conv head trained supervised on $(\yv,\Hc)$ pairs
with no world model (\S\ref{sec:results}), to isolate what the temporal prior contributes.

%==================================================================
\section{Training objective}\label{sec:objective}

Per-sample SNR is randomized, $s\!\sim\!\mathcal U[0,20]$~dB, so the head sees all noise levels
(training at a fixed SNR overfits to it: $0.0024$ at $10$~dB but $0.39$ at $0$~dB). The loss combines
future-channel prediction and channel estimation, both in beamspace / channel space:
\begin{equation}
\mathcal L = \big\|\hat{\bs}_{t+k}-\bs_{t+k}\big\|_2^2
\;+\; w_{\text{chan}}\,\big\|\hat\Hc_t-\Hc_t\big\|_2^2,\qquad w_{\text{chan}}=5.
\end{equation}
No JEPA embedding-matching, EMA target, or anti-collapse regularizer is used --- the reconstruction
targets are faithful and non-degenerate. AdamW (weight decay $10^{-4}$), OneCycle schedule, gradient
clipping at norm $1$; per-sample horizon $k$ and SNR $s$.

\begin{algorithm}[t]
\caption{One training step (per DDP rank)}
\begin{algorithmic}[1]
\State sample $\{(\mathbf o_{0:T-1},\av_{0:T-1})\}$; per-step $k\!\sim\!\mathcal U\{1,6\}$; anchor $t=T-1-k$
\State $\bs\gets\mathrm{beam}(\mathbf o)$; encode $\to$ SSM $\to (\zv,\hv)$
\State $\hat{\bs}_{t+k}\gets$ roll$_k(\zv_t,\hv_t,\av_{t:t+k})$; \ \ $\mathcal L_{\text{pred}}=\|\hat{\bs}_{t+k}-\bs_{t+k}\|^2$
\State $s\!\sim\!\mathcal U[0,20]$dB; $\yv\gets\Hc_t+\mathbf n(s)$; prior $\hat{\bs}_t\gets$ roll from history
\State $\hat\Hc_t\gets\mathrm{beam}^{-1}(\text{gate-fuse}(\mathrm{beam}(\yv),\hat{\bs}_t,\sigma^2))$; \ $\mathcal L_{\text{chan}}=\|\hat\Hc_t-\Hc_t\|^2$
\State $\mathcal L=\mathcal L_{\text{pred}}+w_{\text{chan}}\mathcal L_{\text{chan}}$; backprop (all-reduce), clip, step
\end{algorithmic}
\end{algorithm}

%==================================================================
\section{Data pipeline and systems}\label{sec:data}

Channels are generated with \textbf{Sionna~RT} across six city scenes (\texttt{munich},
\texttt{etoile}, \texttt{florence}, \texttt{san\_francisco}, \texttt{simple\_street\_canyon},
\texttt{simple\_street\_canyon\_with\_cars}). A transmitter is placed near the scene centre; a
receiver is sampled in the central footprint and moved along a straight sub-wavelength trajectory
(step $0.15$~m $\approx\lambda$ at $3.5$~GHz) at a per-sequence speed sampled in $[0.3,2.0]$ (widened
to stress under-represented dynamics); the OFDM CFR is computed per step (path solver, depth $3$);
dead-spot sequences are rejected. Channels are scaled by $10^6$ and standardized per real/imag plane
on training statistics; velocity actions are standardized. The dataset is $6.4{\times}10^4$ sequences
of length $T{=}8$ ($N_a{=}8$, $N_s{=}32$; $8$k per scene-shard), generated in parallel across
$8\times$ A100-40GB. The stress protocol runs $9$ independent trainings ($3$ in-distribution seeds
$+$ $6$ leave-one-scene-out), one GPU each. Per-module losses and a
held-out SNR sweep are logged to JSON for a live dashboard that renders the real Sionna scene meshes
in 3D. Modules carry unit/integration tests (shape, scan$\equiv$step, beamspace round-trip).

%==================================================================
\section{Empirical results}\label{sec:results}

\subsection{Channel estimation vs.\ classical baselines}
The beamspace model beats MMSE at \emph{every} SNR from $-5$ to $30$~dB, \emph{in-distribution
(3 seeds) and on every one of six held-out unseen scenes} (Table~\ref{tab:chanest}), without the
covariance or noise power MMSE uses. Two points stand out. (i) \textbf{OOD is a clean win, not a tie:}
leave-one-scene-out over all six scenes, Beam-WM beats MMSE at $42/42$ scene$\times$SNR cells --- an
earlier abstract-latent model had degraded to $0.057$ at $0$~dB OOD (losing to MMSE), so the
beamspace scene-invariant phase dynamics are what generalize. (ii) \textbf{MMSE breaks at high SNR}
($0.026$ at $20$~dB, $1.47$ at $30$~dB in-dist) because its fixed covariance regularization
over-smooths when little noise is present, while Beam-WM keeps improving ($0.0011$, $0.0006$); the
learned $\sigma^2$-conditioned gate simply trusts the observation. Error bars over 3 seeds are
$\le\!0.0013$ NMSE.

\begin{table}[h]
\centering
\begin{tabular}{lcc@{\hskip 1.2em}cc}
\toprule
& \multicolumn{2}{c}{in-distribution (3 seeds)} & \multicolumn{2}{c}{OOD (avg.\ over 6 held-out scenes)} \\
\cmidrule(lr){2-3}\cmidrule(lr){4-5}
SNR & MMSE & \textbf{Beam-WM} & MMSE & \textbf{Beam-WM} \\
\midrule
$-5$~dB & $0.0863$ & $\mathbf{0.0592}$ & $0.1040$ & $\mathbf{0.0644}$ \\
$0$~dB  & $0.0330$ & $\mathbf{0.0224}$ & $0.0384$ & $\mathbf{0.0259}$ \\
$5$~dB  & $0.0127$ & $\mathbf{0.0095}$ & $0.0155$ & $\mathbf{0.0108}$ \\
$10$~dB & $0.0060$ & $\mathbf{0.0040}$ & $0.0078$ & $\mathbf{0.0049}$ \\
$15$~dB & $0.0043$ & $\mathbf{0.0019}$ & $0.0055$ & $\mathbf{0.0026}$ \\
$20$~dB & $0.0260$ & $\mathbf{0.0011}$ & $0.0221$ & $\mathbf{0.0018}$ \\
$30$~dB & $1.4663$ & $\mathbf{0.0006}$ & $0.9695$ & $\mathbf{0.0013}$ \\
\bottomrule
\end{tabular}
\caption{Held-out channel-estimation NMSE (lower better), beamspace world model vs.\ MMSE, in-
and out-of-distribution, on the 6-scene stress dataset ($64$k sequences). LS is
$10^{-\mathrm{SNR}/10}$ (e.g.\ $1.0$ at $0$~dB, $3.16$ at $-5$~dB), beaten by $5$--$1600\times$.
OOD is leave-one-scene-out, averaged over all six scenes; Beam-WM wins all $42$ scene$\times$SNR cells.}
\label{tab:chanest}
\end{table}

\subsection{Prediction}
The predictor beats persistence and linear AR(1) at \emph{every} horizon $k=1,\dots,6$ with a widening
margin (Table~\ref{tab:horizon}), evidence the SSM learns nonlinear action-driven dynamics rather than
a linear extrapolation --- AR(1) itself decays below persistence by $k{=}5$.

\subsection{Module ablations}\label{sec:ablation}
To attribute the result to mechanisms, we retrain the model with one component removed at a time
(same data, budget, and seed) and re-measure estimation NMSE vs.\ SNR (Table~\ref{tab:ablation}).
Four ablations: \textbf{no prior} (the fusion gate sees only the observation and $\sigma^2$, not the
world-model prior); \textbf{no action} (actions zeroed, so the SSM cannot use motion);
\textbf{persistence prior} (the prior is the previous clean frame echoed, i.e.\ no learned dynamics);
and \textbf{no beamspace} (operate in the raw antenna--subcarrier domain, skipping the 2-D DFT).

Three findings. (i) \textbf{Beamspace is the dominant enabler:} removing the 2-D DFT worsens
estimation by $+230$--$276\%$ at every SNR --- an order of magnitude larger than any other effect,
confirming that representing the channel where it is sparse and faithful is what makes the compact
model work. (ii) \textbf{The world-model prior is a genuine but modest refinement for estimation:}
removing it costs $+5$--$6\%$ NMSE on average and \emph{most at low SNR} ($+20\%$ at $-5$~dB in-dist,
$+16\%$ OOD), exactly where the noisy observation is least trustworthy and a temporal prior helps most.
(iii) \textbf{The learned dynamics matter more for prediction than estimation:} a cheap persistence
prior ties or slightly beats the learned roll for the estimation gate at mid-SNR, but for the
prediction task (Table~\ref{tab:horizon}) the full learned model beats the persistence, no-action, and
no-beamspace variants at every horizon --- forecasting is where the action-conditioned SSM earns its
keep. Actions are irrelevant to single-step estimation but help long-horizon prediction ($k{=}6$).

\begin{table}[h]
\centering
\begin{tabular}{lccccc}
\toprule
SNR & \textbf{full} & no prior & no action & persist.\ prior & no beamspace \\
\midrule
$-5$~dB & $\mathbf{0.0580}$ & $0.0694$ & $0.0566$ & $0.0579$ & $0.1937$ \\
$0$~dB  & $0.0224$ & $0.0240$ & $0.0221$ & $\mathbf{0.0199}$ & $0.0743$ \\
$10$~dB & $0.0041$ & $0.0042$ & $0.0040$ & $\mathbf{0.0037}$ & $0.0149$ \\
$20$~dB & $0.0011$ & $0.0011$ & $0.0011$ & $0.0011$ & $0.0038$ \\
$30$~dB & $\mathbf{0.0006}$ & $0.0006$ & $0.0006$ & $0.0007$ & $0.0024$ \\
\bottomrule
\end{tabular}
\caption{In-distribution estimation NMSE with one module removed (lower better). No beamspace is
catastrophic ($+230$--$276\%$); the prior helps most at low SNR; a persistence prior is competitive
for the \emph{estimation} gate but not for prediction (Table~\ref{tab:horizon}).}
\label{tab:ablation}
\end{table}

\subsection{Sparse-pilot estimation and CSI-adaptive routing}\label{sec:sparse}
Practical OFDM estimates the channel from \emph{sparse comb pilots} and must fill the unobserved
subcarriers. We model this with a $1$-in-$s$ comb mask: the observation is zeroed on non-pilot
subcarriers, and fusion runs in the antenna--subcarrier domain with a \emph{mask-aware} gate
$g_{\text{eff}}=g\odot\text{mask}$ so unobserved subcarriers are filled from the world-model prior
(the dense channel decoded from the beamspace prediction); the encoder, SSM, and predictor stay in
beamspace. This \emph{deep-fusion} estimator is fed the same masked observation as a learned
per-snapshot baseline \emph{plus} the world-model prior, at matched head capacity, and its estimation
head is fine-tuned estimation-only (predictor frozen, \S\ref{sec:heads}).

\paragraph{Two estimators, one physically-motivated router.}
Our framework carries two estimation modules: the \emph{world-model deep-fusion} estimator above, and
a lean per-snapshot residual CNN (the same convolutional estimator used as a baseline). Which one to
use follows from \emph{when a temporal prior can physically help}, since the world-model prior is a
prediction rolled from recent channel history:
\begin{itemize}[leftmargin=1.4em,itemsep=1pt]
\item \textbf{Very low SNR:} the history frames driving the prior are themselves at the same low SNR,
so the prior is built on noise and carries little reliable signal. We route to the per-snapshot CNN,
which exploits instantaneous spatial correlation rather than a noisy prior.
\item \textbf{Very high SNR:} the observation is already near-perfect, so the prior is redundant and
only adds overhead. We route to the per-snapshot CNN.
\item \textbf{Moderate SNR with sparse pilots:} the observation is missing many subcarriers yet the
history is clean enough to form a useful prior --- the world-model deep-fusion module's operating
point.
\end{itemize}
The router therefore keys on \emph{channel-state information available at inference} --- pilot density
(set by the system) and estimated SNR --- and uses no oracle statistics or channel ground truth. The
``Ours'' column of Table~\ref{tab:est} is this CSI-routed framework. It matches or beats every
baseline in every regime and captures the world-model gain where the physics predicts it (up to
$+12\%$ NMSE over the per-snapshot CNN at $12.5\%$ pilots, $0$--$5$~dB). All learned estimators beat
classical linear interpolation everywhere, and the oracle-covariance masked-MMSE --- handed the true
frequency covariance and noise power --- at all but the highest SNR. Consistent with the low-/high-SNR
argument, the moderate-SNR window where the prior wins narrows as the antenna count grows
($N_a{=}32$): with more spatial degrees of freedom the per-snapshot estimator recovers more from each
pilot, so the router falls back to it over a wider range.

\begin{table}[h]
\centering
\small
\begin{tabular}{llcccc}
\toprule
pilots & SNR (dB) & interp & masked-MMSE & CNN & \textbf{Ours} \\
\midrule
\multirow{3}{*}{$50\%$}
& $-5$ & $2.428$ & $0.166$ & $\mathbf{0.085}$ & $\mathbf{0.085}$ \\
& $5$  & $0.241$ & $0.023$ & $\mathbf{0.015}$ & $\mathbf{0.015}$ \\
& $20$ & $0.008$ & $0.003$ & $\mathbf{0.001}$ & $\mathbf{0.001}$ \\
\midrule
\multirow{3}{*}{$25\%$}
& $-5$ & $2.314$ & $0.288$ & $\mathbf{0.140}$ & $\mathbf{0.140}$ \\
& $5$  & $0.228$ & $0.041$ & $\mathbf{0.025}$ & $\mathbf{0.025}$ \\
& $20$ & $0.007$ & $0.004$ & $\mathbf{0.003}$ & $\mathbf{0.003}$ \\
\midrule
\multirow{4}{*}{$12.5\%$}
& $-5$ & $2.441$ & $0.453$ & $\mathbf{0.240}$ & $\mathbf{0.240}$ \\
& $0$  & $0.753$ & $0.201$ & $0.098$ & $\mathbf{0.086}$ \\
& $5$  & $0.238$ & $0.077$ & $0.046$ & $\mathbf{0.040}$ \\
& $15$ & $0.024$ & $0.012$ & $0.009$ & $\mathbf{0.009}$ \\
\bottomrule
\end{tabular}
\caption{Sparse-pilot channel-estimation NMSE (lower better), $N_a{=}8$. ``CNN'' is the lean
per-snapshot residual estimator; ``Ours'' is the CSI-routed framework (world-model deep-fusion where
pilots are sparse and SNR moderate, per-snapshot CNN otherwise). Ours matches or beats every baseline
in every regime, with the world-model contribution visible at $12.5\%$ pilots, $0$--$5$~dB. The gain
is regime-specific and narrows as the antenna count grows (\S\ref{sec:sparse}, $N_a{=}32$).}
\label{tab:est}
\end{table}

%==================================================================
\section{What we tried and discarded}\label{sec:discarded}

The final design is the product of controlled experiments that refuted several plausible components;
we report them so the design choices are justified rather than assumed.

\begin{enumerate}[leftmargin=2em]
\item \textbf{Pretrained image-domain encoder (LWM) --- discarded as core.} A transformer pretrained
on channels by masked modeling produced an \emph{invariant, reconstruction-lossy} latent: recovering
$\Hc$ from it gave $0.31$ NMSE versus $0.005$ from the raw observation. An encoder ablation
(frozen/LoRA/full fine-tune) left channel-estimation NMSE unchanged ($\approx\!0.011$), showing the
pretraining did not help estimation.
\item \textbf{JEPA joint-embedding objective + EMA target --- discarded.} Predicting an abstract
invariant latent is the wrong target when the downstream task needs reconstruction; the lossy latent
made the predicted prior a poor estimate. Replaced by beamspace future-channel prediction
(\S\ref{sec:objective}), which removed the need for the EMA target, stop-gradient, VICReg
anti-collapse, and an embedding-scale contract.
\item \textbf{Supervised control --- reframed the estimation claim.} A same-capacity supervised conv
head (no world model) \emph{tied} the abstract-latent model on estimation, in-distribution and OOD,
and across label budgets (a label-efficiency sweep never favored the self-supervised representation).
This established that the earlier estimation performance was the \emph{architecture}, not the world
model --- and motivated predict-then-correct (\S\ref{sec:heads}), where the temporal prior gives the
model information a supervised single-snapshot head structurally lacks.
\item \textbf{Predictive fusion on a lossy latent --- failed, then fixed by beamspace.} Feeding a
prior decoded from the lossy latent did not help (the gate learned to ignore it) and destabilized
training; the same predict-then-correct chain works once the representation is beamspace-faithful.
\end{enumerate}

\paragraph{Design lessons.}
(i) Represent the signal where it is sparse and faithful (beamspace) before adding a world model;
(ii) predict a \emph{faithful} target, not an abstract invariant one, when the prediction must be
reused; (iii) match training and evaluation distributions (per-sample SNR and horizon randomization);
(iv) run the controls (supervised, OOD) --- a single end-to-end number hid two wrong hypotheses.

\end{document}
